<<<<<<< HEAD
\section{Implications for Research \& Practice}
=======
\section{\MakeUppercase {Implications for Research \& Practice}}
>>>>>>> 5962c492beeeeff7385c8b177888710a79b39489
\label{sec:ImplicationsResearch}
\renewcommand{\baselinestretch}{2.0}\normalsize

<<<<<<< HEAD
The three gaps identified in this review -- conceptual ambiguity, limited empirical grounding, and methodological fragmentation -- do not point to isolated weaknesses. Together, they highlight a wider issue with the way resilience is studied in digitally mediated environments. The issue is not just that the literature is incomplete. Rather, resilience is still too often invoked in general terms, measured inadequately, and only operationalised after disruption has already occurred. The following implications take a different view. They treat resilience as a socio-technical capability that must be defined more precisely, observed in action and linked to effective interventions.\\
\\
Firstly, resilience should not be considered a fixed organisational attribute. It is not merely a residual label attached to organisations that appear to cope well after disruption, either. It is better understood as an emergent socio-technical capability. This is important because the literature often moves too easily between resilience, robustness, recovery and agility\autocite{woods2015four,duchek2020organizational}. The overlap is real, but so are the differences. In digital settings especially, resilience is not confined to one asset, one team, or one managerial disposition. Rather, it emerges through interaction between people, technical systems, routines, escalation pathways and governance arrangements. A good example is cyber incident response, where it isn't about one particular resource but rather how well the analysts, security software, and other reporting protocols, along with other channels for information sharing, can collaborate successfully within time constraints. Resilience framed in such a manner will give IS researchers a better starting point for their analyses, while facilitating a more productive interaction between IS research and dynamic capability studies.\autocite{henfridsson2013generative,boh2023building}, while remaining compatible with systems and ecological traditions \autocite{holling1973resilience,wiener1948cybernetics}.\\
\\
A second implication follows from this almost immediately. Resilience must be made observable. Otherwise, it remains just a slogan. Behavioural and computational indicators are important not because they make resilience appear more scientific, but because they make the concept empirically accountable. Plausible candidates include coordination speed, stakeholder support, community ties, the capacity to reconfigure infrastructures, and variation in response under stress \autocite{mayer2021putting,hollnagel2013resilience,gillard2023efficient}. Once specified, such indicators can be studied through agent-based simulation, network analysis, or probabilistic modelling \autocite{helbing2012agent}. This is especially important in digital spaces, where adaptation is spread out, needs to happen quickly, and is often only partially visible. A simulation model can show how changes in how quickly information is shared or how weak incentives for collaboration can change how people work together during cyber incidents. That is the point. When you can see something, you can compare resilience. It can also be questioned and changed.\\

This also means we spend less time treating compliance artifacts as evidence of resilience. Although audits, certifications, and maturity assessments are not useless, they are weak indicators of adaptive capacity in practice.  \autocite{skopik2016problem}. They inform us of the organizations that claim to be prepared to do so. They divulge scant information regarding the ramifications of infrastructure failure, the dissolution of coordination, or the transformation of circumstances midway. Field-based work is therefore essential. Longitudinal case studies, trace data, and ethnographic approaches are important here. Research in healthcare and public health already shows that resilience can be evident in everyday practice while remaining only partially institutionalized \autocite{kruk2017building,barasa2018what}. The same opportunity exists in digital infrastructures, where incidents and near-misses reveal how organisations improvise, adapt, stall or fail over time. In some cases, disruption may trigger processes of renewal reminiscent of Schumpeter's 'creative destruction', in which routines or infrastructures that no longer fit prevailing conditions are replaced by more adaptive ones  \autocite{schumpeter1942capitalism}. This should not be romanticised. However, it does reinforce the broader point that resilience becomes analytically meaningful when studied in action rather than inferred from documentation.\\
\\
Another implication is that adaptive capacity does not belong to only one level of analysis. Rather, it unfolds across micro-level routines, meso-level infrastructures, macro-level institutions, and inter-organizational ecosystems. This is one reason why single-level explanations often feel inadequate. A cyberattack may begin with analyst decisions at the micro level, spread through infrastructural dependencies at the meso level, and then be redirected or constrained by regulation, liability, or cross-organizational coordination at broader scales. The mechanisms are distributed. IS research is unusually well placed to deal with that because it can combine process tracing, computational modeling, and comparative case analysis without forcing the problem into a single analytical register \autocite{ostrom1990governing,vogus2007organizational,folke2010resilience}. That is a real advantage of the field. It also creates room for indicators -- such as coordination speed, recovery sequence, or adaptive reconfiguration -- that travel across contexts without collapsing into empty abstraction.\\
\\
There is also a practical implication that should be stated plainly. Better concepts and indicators won't matter much if they don't influence how organizations prepare, make decisions, and coordinate. To that end, the field of resilience research must evolve to become more managerially useful without succumbing to managerial clichés. This requires design principles, benchmarks, early warning systems, and governance mechanisms that serve a purpose beyond mere decoration of annual reports \autocite{weick2011managing,lengnick2011developing,calder2021iso}. The risk is obvious: Organizations can demonstrate resilience long before they operationalize it. Metrics theater is not rare. Dashboards exist, escalation paths are documented, and the language of resilience is prevalent. Yet, adaptive capacity remains thin when pressure actually arrives. Therefore, linking measurement to governance is not optional. In practice, this may involve real-time monitoring of the adaptive response, escalation protocols that clarify authority during disruption, or governance arrangements that assign coordination responsibilities before a crisis begins.\\
\\
When considered as a whole, these findings suggest a need for a more rigorous research agenda in the field of IS. The field does not need more generic endorsements of resilience or another round of compliance-friendly abstraction. To achieve this, it is essential to develop more sophisticated concepts, strengthen the empirical foundation, and devise methods that can elucidate the process by which adaptive capacity emerges in the face of ongoing disruption. For scholars, that creates room for more cumulative theory. It also creates room for better evidence. The conversation is shifted for practitioners from symbolic commitment to capabilities that can be observed, tested, and revised. Reframed in these terms, resilience becomes less a rhetorical aspiration than a grounded socio-technical capacity for operating under strain in digitally mediated environments.
=======
Regulated sectors---finance above all---already run on a mature digital industry backbone. Business continuity plans (BCPs) and business continuity management systems aligned with ISO~22301 \autocite{calder2021iso}; business impact analyses; recovery-time and impact-tolerance objectives for important business services under BCBS operational-resilience principles and national regimes such as FINMA guidance and DORA \autocite{bcbs2021operational,european2022regulation,finma2023operational}; software supply-chain dependency mapping \autocite{shen2024understanding}; ICT continuity and response-and-recovery policies; annual backup-and-restore testing; crisis-communication and third-party dependency mapping; and threat-led penetration-testing programs (TIBER-EU, CBEST) \autocite{ecb2018tiber,boe2021cbest}---together constitute state-of-the-art preparedness, not a missing control baseline. On that backbone, organizations add SOCs, Security Information and Event Management (SIEM) telemetry, and sector exercises---then still struggle to show boards how any of it evidences adaptive performance when disruption departs from script (\autocite{rota2024}; Appendix~\ref{appendix:C}). BCPs are invoked; cross-unit handoffs stall; ticket histories and after-action findings rarely alter governance rhythms. Resilience is declared in every committee; enacted adaptation remains hard to demonstrate. The deeper problem is conceptual: resilience is still treated as a validated achievement---maturity attainment, plan certification, audit clearance---when it is more accurately a continuous journey of learning and exploration within and across organizational boundaries. Research and practice must shift from certifying possession of controls to validating that journey itself---whether sensing, reconfiguration, and coordination actually improve as threat and disruption conditions evolve.\\
\\
Section~\ref{sec:RD} turns that reorientation into three research imperatives, scoped in Figure~\ref{fig:enactment_scope} and prioritized in Table~\ref{fig:resilience_interface_opportunity}. \textbf{Define} what counts as adaptive performance---at infrastructure cascade, internal coordination, and cross-firm activation---before BCP templates and audit scores substitute for it \autocite{hillmann2021organizational}. \textbf{Measure} it from digital traces organizations already generate: SOC and SIEM logs, BCP test and invocation records, recovery timelines, coordination metadata, ISAC records. Start with explainable quantitative methods; add machine-learning and explainable-AI layers where volume demands it---orienting commercial analytics platforms toward governance-facing indicators rather than detection tickets that never leave the security function \autocite{francis2014metric,gillard2023efficient}. \textbf{Govern} it: bind those indicators to BCP escalation paths, vendor recovery, board review, and cooperative activation across partners---and to the transient governance that disruption triggers, instrumenting how decision authority migrates and hands back rather than only the steady-state controls---not slide decks assembled after incidents close \autocite{boin2013resilient,skopik2016problem,ecb2018tiber}. Where single-firm logs cannot reveal path-dependent cascade across cloud and supplier dependencies, simulation must supply the counterfactual ensemble \autocite{dawson2011agentbased,kwakkel2016comparing}. IS research sits at the hinge: the traces, infrastructures, and cross-boundary workflows through which resilience is actually produced are already digital.

The implications below translate that reorientation for IS scholarship and for the organizations whose infrastructures, routines, and partnerships now constitute the primary site of adaptive performance under disruption.\\
\\
Firstly, the field must confront an oddly asymmetric blind spot. Management has long theorized and instrumented adaptation to market conditions---gradual repositioning under stable demand and sharp strategic pivot under shock \autocite{teece2007explicating,duchek2019organizational}---yet still frames security and operational resilience chiefly as extensions of risk management and compliance. Continuously adapting, forecasting, and measuring security postures and coordination capabilities in response to evolving threat landscapes remains comparatively under-theorized and under-instrumented. Because resilience is annexed to security and risk management in both research and practice, it largely misses the perspective shift its own logic demands: operating under disruptive conditions is less a control problem than an ongoing business-adaptation problem. Resilience is therefore best understood not as a fixed organisational attribute, nor only as a label assigned to organisations that appear to cope well after disruption, but as an emergent socio-technical capability enacted across those scopes and mechanisms---a journey whose progress, not its declaration, requires validation. The literature still moves between resilience, robustness, recovery, and agility without distinguishing them sharply \autocite{woods2015four}. In digital settings especially, resilience is not confined to one asset, one team, or one managerial disposition; it emerges through interaction between people, technical systems, routines, escalation pathways, and governance arrangements. Cyber incident response illustrates the point: the outcome depends less on any single resource than on how well analysts, security software, reporting protocols, and information-sharing channels collaborate within time constraints. Framed in this way, resilience provides IS researchers a clearer starting point for analysis and a more productive interaction with dynamic capability studies \autocite{henfridsson2013generative,boh2023building} and digital-transformation research \autocite{vial2019understanding,hanelt2021systematic,wessel2021unpacking}, while remaining compatible with systems and ecological traditions \autocite{wiener1948cybernetics,holling1973resilience}.\\
\\
A second implication follows directly. Validating the resilience journey requires making adaptive trajectories observable---not only snapshots of preparedness at a point in time. Behavioural and computational indicators serve that purpose: not by making resilience appear more scientific, but by giving practitioners and researchers a shared evidentiary basis for whether learning and reconfiguration are improving over successive disruptions, exercises, and near-misses. Plausible candidates include coordination speed, stakeholder support, community ties, the capacity to reconfigure infrastructures, and variation in response under stress \autocite{paries2017resilience,gillard2023efficient}. Once specified, such indicators can be studied through agent-based simulation \autocite{helbing2012agent}, network analysis \autocite{panteli2017metrics}, or probabilistic modelling \autocite{kruk2017building}. This is especially valuable in digital spaces, where adaptation is distributed, time-critical, and only partially visible. A simulation model can show, for instance, how variations in the speed of information sharing or the strength of collaboration incentives shape how people coordinate during cyber incidents. Once made visible, the journey can be compared across cases, examined more carefully, and improved deliberately.\\
\\
A further implication cuts across all three imperatives: artificial intelligence is becoming the connective layer through which they can be pursued at scale, in research and in practice. On the measurement and decision side, explainable AI can fuse the heterogeneous signals resilience leaves behind---SOC and SIEM telemetry, coordination metadata, exercise after-action data, and cross-firm sharing records---into decision-ready indicators, provided its reasoning remains legible to operators, boards, and regulators rather than collapsing into opaque risk scores \autocite{berente2021managing,gillard2023efficient}. On the anticipatory side, reinforcement-learning and related generative methods can search large spaces of adversarial and operational contingencies to surface \emph{robust} future scenarios---the regimes a strategy must be able to withstand---whose mastery is precisely what adaptation under deep uncertainty requires \autocite{kwakkel2016comparing,lempert2019robust}. Framed this way, AI does not displace the construct, indicators, and governance developed above; it integrates them, turning dispersed traces into foresight and foresight into rehearsed response. The same opacity that makes such models powerful, however, makes explainability a precondition rather than an enhancement if their outputs are to carry weight in governance.\\

This also implies extending---rather than replacing---the role of compliance artifacts and continuity planning in resilience assessment. Audits, certifications, business continuity plans, maturity assessments, and BCP test records provide essential accountability and have advanced organizational governance considerably; their range can now usefully be complemented by direct measures of adaptive performance \autocite{skopik2016problem,calder2021iso,bcbs2021operational}. They tell us a great deal about how organizations articulate preparedness; combining them with field-based observation captures more of what happens when infrastructure fails, coordination is reorganized, or circumstances shift midway. Longitudinal case studies, trace data, and ethnographic approaches are particularly well suited to that complement. Research in healthcare and public health already shows that resilience is often evident in everyday practice well before it is fully institutionalized \autocite{kruk2017building,barasa2018what}. A similar opportunity exists in digital infrastructures, where incidents and near-misses offer rich material on how organisations improvise, adapt, recover, or fall short over time. In some cases, disruption triggers processes of renewal reminiscent of Schumpeter's 'creative destruction', in which routines or infrastructures that no longer fit prevailing conditions are replaced by more adaptive ones. Without romanticising disruption, this reinforces the broader point that resilience becomes analytically richest when studied in action alongside the documentation that organizations also produce.\\
\\
Another implication is that adaptive capacity does not belong to only one level of analysis. It unfolds across micro-level routines, meso-level infrastructures, macro-level institutions, and inter-organizational ecosystems. This is one reason why single-level explanations capture only part of the picture. A cyberattack may begin with analyst decisions at the micro level, spread through infrastructural dependencies at the meso level, and then be redirected or constrained by regulation, liability, or cross-organizational coordination at broader scales. The mechanisms are distributed. \textcite{bhamra2011resilience}'s synthesis of the resilience literature already pointed to this multi-level architecture, and subsequent work on the transient governance that disruption induces (Section~\ref{sec:RD}), information cultures that convert near-misses into learning \autocite{westrum2004typology}, and polycentric governance of common resources \autocite{ostrom1990governing} provides complementary entry points. IS research is unusually well placed to engage with that complexity because it can combine process tracing, computational modeling, and comparative case analysis without forcing the problem into a single analytical register \autocite{vogus2007organizational,folke2010resilience}. That is a real advantage of the field. It also creates room for indicators that remain anchored in concrete organizational practice as they travel across contexts.\\
\\
There is also a practical implication worth stating plainly. Better concepts and indicators have most value when they shape how organizations prepare, make decisions, and coordinate. The field of resilience research can therefore continue to grow into a more useful body of work, building on the substantial managerial and regulatory effort already devoted to design principles, benchmarks, early warning systems, and governance mechanisms \autocite{calder2021iso,bcbs2021operational,nist2024csf}. A practical caution accompanies this progress: articulated resilience can outpace its operationalization, and many organizations are already aware of the gap. Field interviews \autocite{rota2024} capture this tension directly: boards request dashboards and KPIs, yet defining and resourcing indicators that boards will actually use often remains unresolved (Appendix~\ref{appendix:C}). Dashboards exist, escalation paths are documented, and the language of resilience is widely shared; the next step is to ensure that these structures connect tightly to adaptive capacity when pressure arrives. Linking measurement to governance is therefore central. In practice, this may involve real-time monitoring of the adaptive response \autocite{kuypers2018designing,gillard2023efficient}, escalation protocols that clarify authority during disruption \autocite{westrum2004typology,boin2013resilient}, or governance arrangements that assign coordination responsibilities before a crisis begins \autocite{ostrom1990governing,arnold2004informationsharing}. Some of these arrangements, in the spirit of \textcite{taleb2012antifragile}, can also be designed to generate productive learning from controlled exposure to stressors, as the threat-led testing regimes formalized by CBEST and TIBER-EU illustrate \autocite{ecb2018tiber}.\\
\\
Considered as a whole, these findings invite a more focused resilience research agenda in the field of IS. Building on the conceptual ground already laid by management, organization, and IS scholars, the next step is to refine the constructs, strengthen the empirical foundation, and develop methods that trace how adaptive capacity evolves---not merely whether controls were present---across successive disruptions. For scholars, this creates room for more cumulative theory and richer evidence on the resilience journey. For practitioners, the conversation can shift from claiming preparedness to demonstrating that adaptation is learning, measurable, and governable alongside the strategic commitments organizations have already made. In these terms, resilience moves from an organizing principle to a grounded socio-technical capacity for operating under increased strain in digitally mediated environments \autocite{majchrzak2016designing}.
>>>>>>> 5962c492beeeeff7385c8b177888710a79b39489
